\documentclass[english,twoside, openright, hidelinks, 11 pt, class=report,crop=false]{standalone}
\input{../../preamb}
\input{../bm}

\begin{document}
\footnotesize


\opr{opgfunkligset} \os
\abc{
\item Av \eqref{eks3a} har vi at \vs
\alg{
	x-y&=5 \\
	y&=x-5
}
Av \eqref{eks3b} har vi at \vs
\alg{
	x+y&=9 \\
	y &= 9-x 
}
De to uttrykkene for $ y $ er uttrykket for to rette linjer, som henholdsvis sammenfaller med uttrykkene for $ f(x) $ og $ g(x) $.
\item Når $ f(x)=g(x) $, har vi at
\alg{
x-5 &= 9-x \\
2x&=14 \\
x &= 7
}
Videre er da $ y=9-7=2 $. Altså er $ x=7 $ og $ y=2 $.
}
\newpage
\grubr{opgfunkrectangleandline} \\
For skjæringspunktet til $f$ og $CD$ har vi at
\begin{align*}
	f(x)&=5 \\
	-2x+9&= 5\\
	-2x &= -4 \\
	x &= 2
\end{align*}
Altså er $E=(2, 5)$ den éne skjæringspunktet. For skjæringspunktet til $f$ og $BC$ har vi at $x=4$. Da er $f(4)=-2\cdot4+9=1$. Altså er $F=(4,1)$ det andre skjæringspunktet. $EC=2$ og $FC=4$. Dermed er
\[ A_{\triangle EFC}=\frac{2\cdot4}{2}=4 \]

\fig{opgfunkrectangleandline}

\grubr{opgfunkkvad}
\abc{
\item Ved å gange ut parentesene får vi at
\[ (x+2)(x-4)=x^2-4x+2x-8=x^2-2x-8=x^2-2x-8 \]
Dette tilsvarer funksjonsuttrykket til $f$.
\item Av uttrykket fra a), finner vi at $f=0$ når $x=-2$ og $x=4$.
\item Vi har at
\begin{align*}
	f(-3)&=(-3+2)(-3-4)=(-1)\cdot(-7)=7 \\
	f(5)&=(5+2)(5-4)=7\cdot1 =7
\end{align*}
Altså er $A=(-3, 7)$ og $B=(5, 7)$. For både $A$ og $B$ er horisontalavstanden til bunnpunktet 4.
\item To punkt med lik horisontalavstand til bunnpunktet vil ha samme $y$-verdi.
}

\end{document}